% =============================================================
% บทที่ 1 ภาพรวมของระบบ
% ข้อมูลจาก main repo (README.md, CLAUDE.md)
% =============================================================
\label{ch:overview}

\section{เกี่ยวกับระบบ}

เว็บแอปพลิเคชัน ``ไดโน'' (Dino Khon Kaen) เป็นระบบแนะนำการท่องเที่ยวจังหวัดขอนแก่น ผู้ใช้เข้าใช้งานผ่านเว็บเบราว์เซอร์โดยไม่ต้องติดตั้งโปรแกรมเพิ่ม นักท่องเที่ยวค้นหาสถานที่ ร้านอาหาร คาเฟ่ และกิจกรรมในจังหวัด สอบถามข้อมูลกับน้องไดโนซึ่งเป็นแชตบอต ให้ระบบจัดแผนการเดินทางรายวันให้อัตโนมัติ และสะสมพอยท์ด้วยการสแกน QR Code ณ สถานที่ร่วมรายการเพื่อนำไปแลกของรางวัล ส่วนผู้ดูแลระบบใช้หน้าผู้ดูแลระบบเพิ่มและปรับปรุงข้อมูล ตรวจสอบข้อมูลที่ผู้ใช้แจ้ง และจัดการพอยท์และของรางวัล ขั้นตอนการทำงานหลักของระบบแสดงดังภาพที่~\ref{fig:workflow}

\begin{figure}[H]
\centering
\begin{tikzpicture}[
    node distance=10bp and 42bp,
    box/.style={draw=black, fill=white, rounded corners=3pt, align=center,
      font=\normalfont\small, minimum height=32bp, text width=66bp, inner sep=3bp},
    adm/.style={box, fill=black!12},
    arr/.style={-{Stealth[length=5bp]}, line width=0.8pt, draw=black!60},
    lbl/.style={font=\normalfont\footnotesize, text=black!75, inner sep=2bp}]
  \node[adm] (a) {ผู้ดูแลระบบเพิ่ม\\และปรับปรุงข้อมูล};
  \node[box, right=of a] (b) {ค้นหาสถานที่\\และกิจกรรม};
  \node[box, right=of b] (c) {ถามแชตบอตและ\\วางแผนทริปด้วย AI};
  \node[box, right=of c] (d) {เดินทาง\\ตามแผน};
  \node[box, below=34bp of d] (e) {สแกน QR Code\\สะสมพอยท์};
  \node[adm, left=of e] (f) {แลกของรางวัล\\ที่เคาน์เตอร์};
  \node[box, below=34bp of b] (g) {แจ้งข้อมูล\\ไม่ถูกต้อง};
  \draw[arr] (a) -- node[lbl, above]{เผยแพร่} (b);
  \draw[arr] (b) -- (c);
  \draw[arr] (c) -- (d);
  \draw[arr] (d) -- (e);
  \draw[arr] (e) -- (f);
  \draw[arr] (b) -- node[lbl, right]{พบข้อมูลไม่ตรง} (g);
  \draw[arr] (g.west) -| node[lbl, pos=0.3, below]{เข้าคิวตรวจสอบ} (a.south);
\end{tikzpicture}
\caption{ขั้นตอนการทำงานหลักของระบบ (กรอบพื้นสีเทาคือขั้นตอนของผู้ดูแลระบบ)}
\label{fig:workflow}
\end{figure}

\section{ผู้ใช้งานและสิทธิ์การใช้งาน}
\label{sec:roles}

ระบบแบ่งผู้ใช้งานเป็น 3 บทบาท ได้แก่ ผู้ใช้งานทั่วไป (Guest) ซึ่งเข้าชมเว็บไซต์โดยไม่ได้เข้าสู่ระบบ สมาชิก (Tourist) ซึ่งสมัครสมาชิกและเข้าสู่ระบบแล้ว และผู้ดูแลระบบ (Admin) แต่ละบทบาทใช้งานฟังก์ชันได้ตามตารางที่~\ref{tab:roles} เครื่องหมาย \cmark{} หมายถึงใช้งานได้ และ -- หมายถึงไม่มีสิทธิ์

\begin{table}[H]
\caption{สิทธิ์การใช้งานของผู้ใช้แต่ละบทบาท}
\label{tab:roles}
\centering
\begin{tabular}{|L{\dimexpr\linewidth-5.7cm-8\tabcolsep-5\arrayrulewidth\relax}|C{1.9cm}|C{1.9cm}|C{1.9cm}|}
\hline
\textbf{ฟังก์ชัน} & \textbf{ผู้ใช้งานทั่วไป} & \textbf{สมาชิก} & \textbf{ผู้ดูแลระบบ} \\ \hline
ค้นหาและดูรายละเอียดสถานที่และกิจกรรม & \cmark & \cmark & -- \\ \hline
สอบถามข้อมูลกับแชตบอต & \cmark & \cmark & -- \\ \hline
วางแผนทริปด้วย AI & \cmark & \cmark & -- \\ \hline
ตอบแบบสอบถามความพึงพอใจ & \cmark & \cmark & -- \\ \hline
บันทึกและจัดการแผนในเมนูทริปของฉัน & -- & \cmark & -- \\ \hline
สแกน QR Code สะสมพอยท์และแลกของรางวัล & -- & \cmark & -- \\ \hline
แจ้งข้อมูลไม่ถูกต้องและเสนอกิจกรรมใหม่ & -- & \cmark & -- \\ \hline
แก้ไขโปรไฟล์ เปลี่ยนรหัสผ่านหรืออีเมล และลบบัญชี & -- & \cmark & -- \\ \hline
จัดการสถานที่ กิจกรรม และฐานความรู้แชตบอต & -- & -- & \cmark \\ \hline
ตรวจคิวรายงานข้อมูล คำขอกิจกรรม และสถานที่จาก Google & -- & -- & \cmark \\ \hline
จัดการ QR Code ของรางวัล และแลกของรางวัลที่เคาน์เตอร์ & -- & -- & \cmark \\ \hline
จัดการบัญชีผู้ใช้ ดูทริปของผู้ใช้ และดูผลความพึงพอใจ & -- & -- & \cmark \\ \hline
\end{tabular}
\end{table}

บัญชีผู้ดูแลระบบสร้างด้วยสคริปต์ตามหัวข้อ~\ref{sec:create-admin} และสมัครผ่านหน้าเว็บไม่ได้ ผู้ดูแลระบบยังเปิดหน้าเว็บสาธารณะได้เช่นเดียวกับผู้ใช้งานทั่วไป

\section{ข้อจำกัดของระบบ}

\begin{enumerate}
\item ข้อมูลสถานที่ครอบคลุมเฉพาะจังหวัดขอนแก่น โดยข้อมูลหลักมาจาก Google Places API และข้อมูลที่ผู้ดูแลระบบเพิ่มเอง จึงอาจไม่ตรงกับความเป็นจริงในช่วงที่ยังไม่ได้ซิงก์ข้อมูลรอบถัดไป
\item แชตบอตและระบบวางแผนทริปใช้โมเดลภาษาขนาดใหญ่ แม้จำกัดการตอบให้อ้างอิงข้อมูลในระบบ แต่ยังมีโอกาสตอบคลาดเคลื่อน ผู้ใช้ควรตรวจสอบข้อมูลสำคัญก่อนเดินทาง
\item QR Code แต่ละจุดให้พอยท์ได้ครั้งเดียวต่อบัญชี และการแลกของรางวัลต้องทำที่เคาน์เตอร์โดยเจ้าหน้าที่ ระบบยังไม่มีการแลกออนไลน์
\item ระบบเป็นเว็บแอปพลิเคชันที่ใช้ผ่านเบราว์เซอร์ ไม่มีแอปพลิเคชันบนมือถือแยก และการสแกน QR Code ต้องอนุญาตให้เบราว์เซอร์ใช้กล้องและตำแหน่งที่ตั้ง
\end{enumerate}

\section{การใช้คู่มือ}
\label{sec:howto-read}

คู่มือแบ่งเนื้อหาตามกลุ่มผู้อ่าน ผู้อ่านเปิดเฉพาะบทที่เกี่ยวข้องกับงานของตนได้ตามตารางที่~\ref{tab:readers}

\begin{table}[H]
\caption{บทที่ผู้อ่านแต่ละกลุ่มควรอ่าน}
\label{tab:readers}
\centering
\begin{tabular}{|L{4cm}|L{\dimexpr\linewidth-8cm-6\tabcolsep-4\arrayrulewidth\relax}|L{4cm}|}
\hline
\textbf{ผู้อ่าน} & \textbf{งานที่ทำ} & \textbf{บทที่ควรอ่าน} \\ \hline
นักท่องเที่ยว & ค้นหาข้อมูล วางแผนทริป และสะสมพอยท์
  & บทที่~\ref{ch:start}, \ref{ch:tourist} และ~\ref{ch:trouble} \\ \hline
เจ้าหน้าที่ผู้ดูแลระบบ & จัดการข้อมูล ตรวจคิวรายงาน และแลกของรางวัลผ่านหน้าเว็บ
  & บทที่~\ref{ch:start}, \ref{ch:admin-web} และ~\ref{ch:trouble} \\ \hline
ผู้ดูแลเครื่องแม่ข่าย & ติดตั้ง ตั้งค่า ตรวจสถานะ และปรับรุ่นระบบ
  & บทที่~\ref{ch:install}, \ref{ch:admin} และ~\ref{ch:trouble} \\ \hline
นักพัฒนา & แก้ไขและพัฒนาระบบต่อ
  & บทที่~\ref{ch:install}, \ref{ch:dev} และภาคผนวก~ข \\ \hline
\end{tabular}
\end{table}

\section{สัญลักษณ์และรูปแบบที่ใช้ในคู่มือ}

คู่มือนี้ใช้รูปแบบข้อความตามตารางที่~\ref{tab:conventions} เพื่อให้ผู้อ่านแยกองค์ประกอบบนหน้าจอออกจากเนื้อความได้ง่าย

\begin{table}[H]
\caption{สัญลักษณ์และรูปแบบข้อความที่ใช้ในคู่มือ}
\label{tab:conventions}
\centering
\begin{tabular}{|C{4.4cm}|L{8.8cm}|}
\hline
\textbf{รูปแบบ} & \textbf{ความหมาย} \\ \hline
``เข้าสู่ระบบ'' & ชื่อปุ่ม เมนู แท็บ หรือช่องกรอกตามที่ปรากฏบนหน้าจอ \\ \hline
``โปรไฟล์ของฉัน > ความปลอดภัย'' & เมนูและแท็บที่ต้องเลือกตามลำดับ \\ \hline
Ctrl+F5 & แป้นพิมพ์ที่ต้องกดพร้อมกัน \\ \hline
หมายเลขสีแดงบนภาพ & หมายเลขกำกับบนภาพหน้าจอ อธิบายเป็นข้อตามลำดับใต้ภาพ \\ \hline
\cmd{npm run dev} & คำสั่งหรือข้อความที่ต้องพิมพ์ตามตัวอักษร \\ \hline
\end{tabular}
\end{table}
